% Generated by code/make_tables.py from results/metrics.json.
% Do not edit by hand: re-run the pipeline instead.
\begin{table}[htbp]
    \centering
    \caption{Training accuracy against test accuracy under the
    leave-one-specimen-out protocol. Each fold's model is scored on the
    micrographs it was trained on and on the specimen it was held out from, by
    the same code and against the same annotations. The gap is train minus test
    in each metric's own units, so the model favours what it has already seen
    when the $\mathrm{AP}_{50}$ gap is positive and when the counting gap is
    negative. The training figure is not a result in
    its own right -- it is measured on data the optimiser was given -- but the
    distance between the two columns is: it separates a model that has memorised
    its training specimens from one that is limited by the difficulty of the
    task. No pooled figure is given because a micrograph belongs to the training
    partition of three of the four folds, so pooling would weight specimens by
    how often they recur.}
    \label{tab:generalisation}
    \begin{tabular}{lrrrrrr}
        \toprule
        & \multicolumn{3}{c}{$\mathrm{AP}_{50}$} &
        \multicolumn{3}{c}{Counting error (\%)} \\
        \cmidrule(lr){2-4} \cmidrule(lr){5-7}
        Method & Train & Test & Gap & Train & Test & Gap \\
        \midrule
        Mask R-CNN               & 0.437 & 0.298 & +0.139 & 32.6 & 41.1 & -8.5 \\
        Faster R-CNN             & 0.412 & 0.305 & +0.107 & 35.7 & 37.8 & -2.2 \\
        Mask R-CNN, $3\times$ bicubic & 0.459 & 0.373 & +0.086 & 46.0 & 42.7 & +3.3 \\
        \bottomrule
    \end{tabular}
\end{table}
